\section{术语消化：本期关键词索引}

\begin{longtable}{p{0.25\textwidth}p{0.34\textwidth}p{0.32\textwidth}}
\toprule
术语 & 一句话解释 & 在本期中的作用 \\
\midrule
Full Attention & 标准 Softmax 全局注意力 & 表达强，但长上下文成本高。 \\
Linear Attention & 近似线性复杂度的注意力变体 & Kimi/Qwen 等路线重点探索。 \\
Sparse Attention & 只选择部分 token 参与注意力 & DeepSeek 路线重点探索。 \\
KV Cache & 推理时保存历史 key/value 的缓存 & 长上下文显存瓶颈。 \\
KDA & Kimi Delta Attention & Kimi Linear 的核心线性注意力模块。 \\
Delta Rule & 一类递推更新规则 & KDA/Gated DeltaNet 的历史线索。 \\
Gated DeltaNet & 带门控的 DeltaNet 变体 & KDA 的基础之一。 \\
Decay & 记忆衰减率 & 细粒度 decay 改善 hidden state 使用。 \\
Hybrid Attention & Full 与 Linear/Sparse 混合架构 & 当前实用路线。 \\
Scaling Ladder & 逐级放大规模验证架构 & Kimi 内部评估机制。 \\
MoE & Mixture of Experts，混合专家 & FFN 侧已被证明的架构突破。 \\
FLOPs & 浮点运算量 & 比较计算成本与效率。 \\
NoPE / RoPE & 无位置编码 / 旋转位置编码 & 训练细节与位置建模相关。 \\
Chunkwise Algorithm & 分块并行算法 & 让 recurrence 类算法能并行训练。 \\
Hardware-friendly & 硬件亲和 & 算法能否大规模落地的关键。 \\
\bottomrule
\end{longtable}

\subsection{本章小结}

本期术语密度很高，但核心关系很清楚：长上下文让 Full Attention 成本上升，Linear/Sparse/Hybrid 都在尝试降低成本；真正有前途的路线必须同时兼顾表达能力、效率和硬件亲和。

